\chapter{CÁC GIẢI PHÁP VÀ ĐÓNG GÓP NỔI BẬT}

Chương này trình bày ba đóng góp kỹ thuật nổi bật của hệ thống, mỗi đóng góp được trình bày
theo cấu trúc: (1)~đặt vấn đề, (2)~giải pháp được đề xuất, (3)~kết quả đạt được.

Agent:

Trong khuôn khổ hệ thống đã được thiết kế và xây dựng tại chương 4, ba đóng góp kỹ thuật sau đây được đánh giá là có tính mới và ứng dụng thực tiễn cao nhất, đồng thời phản ánh các quyết định kiến trúc quan trọng trong quá trình phát triển: (i)~pipeline phân loại và dịch tài liệu đa định dạng với thuật toán nhận diện PDF thông minh trong mục 5.1, (ii)~hệ thống quản lý thuật ngữ ba cấp với quy trình đề xuất cộng đồng theo ngưỡng trong mục 5.2, và (iii)~cơ chế inject ngữ cảnh bảng thuật ngữ qua thread-local storage trong mục 5.3.

% =========================================================
\section{Pipeline dịch tài liệu đa định dạng với phân loại PDF thông minh}

\subsection{Đặt vấn đề}

PDF là định dạng tài liệu phổ biến nhất trong môi trường doanh nghiệp, song đây cũng là định
dạng phức tạp nhất để dịch thuật vì PDF tồn tại dưới hai dạng hoàn toàn khác nhau:

\begin{itemize}
  \item \textbf{PDF native-text (có thể chỉnh sửa):} PDF được tạo trực tiếp từ phần mềm văn
    phòng (Word, Excel) hoặc in ra PDF từ ứng dụng. Loại này chứa các đối tượng văn bản
    (\texttt{stream}) có thể trích xuất trực tiếp.
  \item \textbf{PDF scan (dạng ảnh):} PDF được tạo bằng cách scan tài liệu giấy hoặc chụp ảnh.
    Toàn bộ nội dung trang là ảnh bitmap, không chứa đối tượng văn bản có thể trích xuất.
  \item \textbf{PDF OCR overlay:} PDF scan đã được chạy qua phần mềm OCR và nhúng thêm một lớp
    văn bản không hiển thị phía trên ảnh. Đây là trường hợp gây nhầm lẫn nhất vì có văn bản
    nhưng không phải native.
\end{itemize}

Nếu áp dụng sai pipeline (ví dụ: cố gắng trích xuất văn bản từ PDF scan), kết quả sẽ trống hoặc
sai hoàn toàn. Ngược lại, nếu dùng OCR cho PDF native-text, sẽ tốn thời gian và giảm chất lượng
không cần thiết.

\subsection{Giải pháp: Thuật toán phân loại PDF đa tiêu chí}

Đồ án đề xuất hàm \texttt{is\_pdf\_truly\_editable(pdf\_path)} phân tích từng trang PDF theo
ba tiêu chí kết hợp để phân loại chính xác:

\begin{enumerate}
  \item \textbf{Tiêu chí 1 -- Văn bản native:} Trích xuất văn bản từ trang bằng
    \texttt{pypdf.PdfReader.extract\_text()}. Nếu có $\geq 20$ ký tự, trang được coi là có
    nội dung văn bản.
  \item \textbf{Tiêu chí 2 -- Phát hiện ảnh toàn trang:} Kiểm tra \texttt{/XObject} trong
    \texttt{/Resources} của trang. Nếu có ảnh (\texttt{/Subtype = /Image}) có kích thước pixel
    lớn hơn ngưỡng tỉ lệ so với kích thước trang (tỉ lệ $\geq 2.0$), trang được phân loại là
    ảnh toàn trang (scan).
  \item \textbf{Tiêu chí 3 -- Metadata OCR:} Kiểm tra \texttt{/Producer}, \texttt{/Creator}
    trong metadata PDF; nếu chứa các từ khóa OCR (``tesseract'', ``adobe pdf output intent''),
    trang được đánh dấu là OCR overlay.
\end{enumerate}

Kết quả phân loại từng trang được tổng hợp theo logic:

\begin{lstlisting}[language=Python, caption={Logic phân loại PDF (translate.py)}]
# Phan loai moi trang:
if has_text and not has_full_image:
    native_text_pages += 1       # van ban nguyen ban
elif has_text and has_full_image:
    ocr_overlay_pages += 1       # scan + OCR overlay
elif (not has_text) and has_full_image:
    image_only_pages += 1        # scan khong co OCR

# Quyet dinh cuoi cung:
if native_text_pages > 0:
    return True   # PDF co the chinh sua -> pipeline ConvertAPI->DOCX
# Khong co trang native text:
if ocr_overlay_pages > 0 or image_only_pages > 0:
    return False  # PDF scan -> pipeline OCR Azure
return False      # Khong xac dinh -> chuyen ve OCR de an toan
\end{lstlisting}

Sau khi phân loại, hệ thống định tuyến đến pipeline phù hợp:

Agent:

\begin{figure}[H]
\centering
\fbox{\parbox{0.9\textwidth}{\centering\vspace{0.5cm}
\textit{[Hình: Sơ đồ quyết định pipeline PDF]}\\[0.3cm]
\textbf{Image Description:} Flowchart: Input PDF → \texttt{is\_pdf\_truly\_editable()} →
\newline
Nhánh TRUE: ConvertAPI (PDF→DOCX) → \texttt{translate\_docx()} → Output DOCX dịch.
\newline
Nhánh FALSE: Azure Document Intelligence OCR (PDF→text+layout) →
\texttt{translate\_text\_with\_glossary()} → Vẽ lại PDF với text dịch → Output PDF dịch.
\vspace{0.5cm}}}
\caption{Sơ đồ quyết định pipeline dịch PDF}
\label{fig:pdf_pipeline}
\end{figure}

Agent:
\begin{lstlisting}[language=bash, caption={Mermaid source -- Pipeline phân loại và dịch PDF}]
flowchart TD
    INPUT["Input: Tệp PDF"] --> CHECK{"is_pdf_truly_editable?\n(3 tiêu chí:\n1. has extractable text\n2. NOT all images\n3. char/word ratio OK)"}

    CHECK -->|"TRUE: PDF native text"| CONV["ConvertAPI\nPDF -> DOCX"]
    CHECK -->|"FALSE: PDF scan / image"| OCR["Azure Document Intelligence\nOCR Pipeline"]

    CONV --> DOCX_TRANS["translate_docx()\nBatch XML paragraph translation\n+ glossary injection"]
    DOCX_TRANS --> OUT_DOCX["Output: DOCX da dich\n(bao toan dinh dang)"]

    OCR --> EXTRACT["Trich xuat text + layout\n(bang, doan van, tieu de)"]
    EXTRACT --> TRANS["translate_text_with_glossary()\nper segment -- thread-local context"]
    TRANS --> DRAW["Ve lai PDF\nvoi text da dich + font embedding"]
    DRAW --> OUT_PDF["Output: PDF da dich"]
\end{lstlisting}

\subsection{Kết quả đạt được}

\begin{itemize}
  \item Hệ thống phân loại chính xác PDF native-text, PDF scan và PDF OCR overlay trong thử
    nghiệm với bộ tài liệu kỹ thuật thực tế (báo cáo chất lượng, hướng dẫn vận hành).
  \item PDF native-text được xử lý qua ConvertAPI→DOCX, giữ nguyên bố cục trang, bảng biểu và
    định dạng chữ sau khi dịch.
  \item PDF scan được xử lý qua Azure Document Intelligence, nhận diện đúng bố cục bảng và
    văn bản nhiều cột với độ chính xác OCR cao hơn so với Tesseract trong thử nghiệm với
    tài liệu tiếng Nhật và tiếng Trung.
  \item Pipeline xử lý lỗi theo từng ngôn ngữ đích độc lập: nếu một ngôn ngữ thất bại, các
    ngôn ngữ còn lại vẫn trả kết quả.
\end{itemize}

% =========================================================
\section{Hệ thống quản lý thuật ngữ ba cấp với quy trình đề xuất cộng đồng}

\subsection{Đặt vấn đề}

Thuật ngữ kỹ thuật chuyên ngành trong môi trường doanh nghiệp (tên sản phẩm, tên quy trình
sản xuất, ký hiệu tiêu chuẩn nội bộ) thường không có trong từ điển mặc định của Google Translation API. Nếu không có cơ chế quản lý thuật
ngữ, cùng một khái niệm kỹ thuật sẽ được dịch khác nhau ở mỗi lần dịch, gây nhầm lẫn trong
vận hành.

Yêu cầu đặt ra:
\begin{itemize}
  \item Cần có kho từ điển thuật ngữ dùng chung toàn công ty (Common Library).
  \item Thuật ngữ phải do người dùng thực tế đóng góp (crowdsourced), không do admin nhập thủ
    công hoàn toàn.
  \item Cần cơ chế kiểm soát chất lượng: thuật ngữ phải được duyệt trước khi áp dụng toàn hệ
    thống.
  \item Người dùng cũng cần có thư viện cá nhân để lưu thuật ngữ riêng chưa chia sẻ.
\end{itemize}

\subsection{Giải pháp: Kiến trúc thuật ngữ ba cấp}

Đồ án thiết kế hệ thống quản lý thuật ngữ theo ba cấp độ:

\textbf{Cấp 1 -- Thư viện cá nhân (Private Library):} Mỗi người dùng có bảng
\texttt{PrivateKeyword} riêng, lưu từ điển đa ngôn ngữ qua trường \texttt{translations}
(JSONField, dạng \texttt{\{"vi": "...", "ja": "...", ...\}}) và trường \texttt{note}. Chỉ
người dùng đó nhìn thấy và có thể sử dụng khi dịch với \texttt{library\_mode=private}.

\textbf{Cấp 2 -- Hàng chờ đề xuất (Suggestion Queue):} Khi người dùng muốn chia sẻ thuật
ngữ lên thư viện chung, hệ thống lưu vào \texttt{KeywordQueue}. Hệ thống đếm số người dùng
khác nhau đề xuất cùng thuật ngữ; khi đạt ngưỡng \texttt{min\_suggesters\_for\_queue} (mặc
định~2), hệ thống tự động tạo \texttt{KeywordSuggestion} với \texttt{status=pending} và gửi
thông báo đến Library Keeper.

\textbf{Cấp 3 -- Thư viện chung (Common Library):} Library Keeper/Admin duyệt các đề xuất.
Khi duyệt (\texttt{status=approved}), hệ thống trigger background thread:
\begin{enumerate}
  \item Tạo tệp CSV từ tất cả \texttt{KeywordSuggestion} đã duyệt.
  \item Upload CSV lên Google Cloud Storage.
  \item Gọi \texttt{manage\_glossary(mode=1)} để cập nhật 45 glossary trên Google Translation
    API (toàn bộ C(n,2) cặp với $n$ ngôn ngữ active, lấy động qua
    \texttt{generate\_pairs\_from\_db()}).
\end{enumerate}

Agent:

\begin{figure}[H]
\centering
\fbox{\parbox{0.9\textwidth}{\centering\vspace{0.5cm}
\textit{[Hình: Sơ đồ hệ thống thuật ngữ ba cấp]}\\[0.3cm]
\textbf{Image Description:} Sơ đồ ba cấp:
\newline
Cấp 1 (Private Library): PrivateKeyword table ←→ User (CRUD). Khi dịch với mode=private:
thread-local context → inject private glossary vào Translation API.
\newline
Cấp 2 (Suggestion Queue): User nhấn ``Đề xuất'' → KeywordQueue → Hệ thống đếm unique
users → Nếu $\geq$ ngưỡng → tạo KeywordSuggestion(pending) → notify Library Keeper.
\newline
Cấp 3 (Common Library): Library Keeper duyệt → status=approved → background thread →
CSV → GCS upload → Google Translation API update 45 glossaries (C(10,2)). Khi dịch với mode=common:
thread-local context → inject common glossary vào Translation API.
\vspace{0.5cm}}}
\caption{Kiến trúc hệ thống quản lý thuật ngữ ba cấp}
\label{fig:library_arch}
\end{figure}

Agent:
\begin{lstlisting}[language=bash, caption={Mermaid source -- Kien truc he thong quan ly thuat ngu ba cap}]
flowchart TB
    subgraph L1["Cap 1: Private Library"]
        direction LR
        USER["Nguoi dung"] -->|"CRUD thuat ngu"| PK[("PrivateKeyword\ntranslations: JSONField\n{vi:.., ja:.., en:..}")]
    end

    subgraph L2["Cap 2: Suggestion Queue"]
        direction TB
        PK -->|"De xuat"| KQ[("KeywordQueue\nis_processed: False")]
        KQ -->|"Background job:\ndem unique users >= min_suggesters"| KS[("KeywordSuggestion\nstatus=pending")]
        KS -->|"Thong bao"| LK["Library Keeper"]
    end

    subgraph L3["Cap 3: Common Library"]
        direction TB
        LK -->|"Duyet / Tu choi"| KS_A[("KeywordSuggestion\nstatus=approved")]
        KS_A -->|"async_manage_common_glossaries()\nbackground thread"| CSV["CSV file\n(pair rows)"]
        CSV -->|"upload_csv_to_gcs()"| GCS[("Google Cloud Storage\nglossary CSV")]
        GCS -->|"manage_glossary()\ncreate / update"| GCTA["Google Translation API\n45 glossary pairs (C(10,2))"]
    end

    subgraph TRANS["Qua trinh dich"]
        direction LR
        MODE{"library_mode?"}
        MODE -->|"mode=private"| PGLOSS["Private Glossary\n_user_ID"]
        MODE -->|"mode=common"| CGLOSS["Common Glossary\ntoray_translation_glossary_*"]
    end

    GCTA --> CGLOSS
    PK -.->|"private mode"| PGLOSS
\end{lstlisting}

Cơ chế ngưỡng đề xuất giúp lọc các thuật ngữ được ít người quan tâm và đảm bảo chỉ những thuật
ngữ được nhiều người dùng xác nhận mới vào hàng chờ duyệt, giảm tải cho Library Keeper.

\subsection{Kết quả đạt được}

\begin{itemize}
  \item Hệ thống thuật ngữ hoạt động tự động từ đầu đến cuối: từ đề xuất của người dùng đến
    áp dụng vào Translation API mà không cần thao tác thủ công trên Google Console.
  \item 45 cặp glossary (C(10,2) với 10 ngôn ngữ mặc định, tự động mở rộng theo số ngôn ngữ
    active trong DB) được đồng bộ tự động sau mỗi lần duyệt thuật ngữ, bao phủ toàn bộ các
    cặp ngôn ngữ đang kích hoạt trong hệ thống.
  \item Cơ chế cập nhật bất đồng bộ (background thread) đảm bảo API dịch thuật không bị chờ
    đợi trong quá trình cập nhật glossary (có thể mất 30--180~giây cho mỗi glossary).
  \item Người dùng nhận thông báo in-app khi đề xuất được duyệt hoặc từ chối.
\end{itemize}

% =========================================================
\section{Cơ chế inject ngữ cảnh glossary qua Thread-Local Storage}

\subsection{Đặt vấn đề}

Module dịch thuật (\texttt{translate\_docx}, \texttt{translate\_xlsx}, \texttt{translate\_pptx},
v.v.) được thiết kế độc lập, với chữ ký hàm nhận vào \texttt{(file\_path, target\_lang,
source\_lang)}. Chúng không có thông tin về ngữ cảnh HTTP request: user là ai, chế độ thư viện
được chọn là gì.

Khi thêm tính năng chọn chế độ thư viện (\texttt{library\_mode=none/common/private}), nảy
sinh thách thức: làm thế nào để truyền thông tin \texttt{(library\_mode, user\_id)} qua nhiều
lớp hàm mà không phải thay đổi chữ ký hàm của tất cả các module dịch?

Thay đổi chữ ký hàm sẽ dẫn đến: phải sửa tất cả module, phá vỡ tính độc lập của từng module,
và tạo sự phụ thuộc giữa các module dịch thuật với tầng API.

\subsection{Giải pháp: Thread-Local Storage}

Đồ án sử dụng \textbf{thread-local storage} (biến cục bộ theo luồng) để inject ngữ cảnh
glossary mà không cần thay đổi chữ ký hàm của bất kỳ module nào.

\begin{lstlisting}[language=Python, caption={Thread-local context trong translate\_text.py}]
import threading
_translation_context = threading.local()

def set_translation_context(library_mode: str, user_id: int):
    _translation_context.library_mode = library_mode
    _translation_context.user_id = user_id

def clear_translation_context():
    _translation_context.library_mode = None
    _translation_context.user_id = None

def translate_text_with_glossary(text, glossary_id,
                                  source_lang_code, target_lang_code):
    # Doc nguon context tu thread-local
    library_mode = getattr(_translation_context,
                           'library_mode', 'common')
    user_id = getattr(_translation_context, 'user_id', None)

    # Chon glossary_id dua tren mode va user
    if library_mode == 'none':
        effective_glossary_id = None
    elif library_mode == 'private' and user_id:
        effective_glossary_id = (
            f"{glossary_id}_{user_id}"  # private glossary
        )
    else:  # 'common'
        effective_glossary_id = glossary_id

    # Goi Google Translation API voi glossary_id da xac dinh
    return _call_translation_api(text, effective_glossary_id,
                                  source_lang_code, target_lang_code)
\end{lstlisting}

Trong \texttt{TranslateFileView.post()}:

\begin{lstlisting}[language=Python, caption={Thiết lập và xóa context (translate.py -- TranslateFileView)}]
# Truoc khi dich:
set_translation_context(
    library_mode=library_mode,  # 'none'/'common'/'private'
    user_id=request.user.id
)
# Toan bo cac module dich (translate_docx, translate_xlsx...)
# tu dong doc context tu thread-local, khong can truyen tham so
for lang in target_languages:
    result = translate_document(temp_path, lang, ...)

# Sau khi dich (trong finally):
clear_translation_context()
\end{lstlisting}

\textbf{Tại sao thread-local là giải pháp phù hợp?}

Mỗi HTTP request trong Django được xử lý trong một luồng riêng. Thread-local đảm bảo ngữ cảnh
của request này không can thiệp vào request khác đang chạy đồng thời, điều mà biến toàn cục
không thể đảm bảo. Đây cũng chính xác là trường hợp sử dụng lý tưởng của thread-local: truyền
thông tin ngữ cảnh trong phạm vi một luồng xử lý mà không cần truyền tường minh qua từng hàm.

Một điểm quan trọng được ghi chú trong source code: module \texttt{translate\_text} phải được
import theo tên trực tiếp (\texttt{import translate\_text}), không phải qua
\texttt{Translate\_v2.translate\_text}, để đảm bảo cả \texttt{TranslateFileView} và các module
dịch chia sẻ \textbf{cùng một instance module} và do đó cùng một đối tượng
\texttt{\_translation\_context}.

\subsection{Kết quả đạt được}

\begin{itemize}
  \item Toàn bộ các module dịch thuật (DOCX, XLSX, PPTX, PDF, text) đều sử dụng ngữ cảnh
    glossary đúng mà không cần thay đổi chữ ký hàm.
  \item Chế độ thư viện được áp dụng nhất quán qua tất cả các định dạng tài liệu.
  \item Hệ thống an toàn với nhiều request đồng thời: thread-local ngăn xung đột ngữ cảnh
    giữa các request.
  \item Dễ mở rộng: để thêm thông tin ngữ cảnh mới (ví dụ: ngôn ngữ ưu tiên của user), chỉ
    cần thêm field vào \texttt{\_translation\_context} mà không ảnh hưởng module nào khác.
\end{itemize}

Agent:

Kết chương 5. Chương 5 đã trình bày chi tiết ba đóng góp kỹ thuật nổi bật của đồ án. Thứ nhất, thuật toán phân loại PDF đa tiêu chí \texttt{is\_pdf\_truly\_editable()} giải quyết vấn đề định tuyến pipeline cho PDF native-text và PDF scan mà không cần can thiệp thủ công, đảm bảo chất lượng dịch tối ưu cho từng loại tài liệu. Thứ hai, kiến trúc thuật ngữ ba cấp --- Private Library, Suggestion Queue và Common Library --- cho phép toàn bộ quy trình từ đề xuất cộng đồng đến cập nhật 45 cặp glossary (C(10,2) với 10 ngôn ngữ mặc định) diễn ra tự động, không yêu cầu thao tác thủ công trên Google Cloud Console. Thứ ba, cơ chế thread-local context injection cho phép tích hợp chế độ thư viện vào tất cả module dịch thuật mà không phá vỡ tính độc lập và khả năng tái sử dụng của từng module, đồng thời đảm bảo an toàn khi xử lý nhiều request đồng thời. Ba đóng góp này kết hợp với nhau tạo nên một hệ thống dịch thuật tài liệu doanh nghiệp hoàn chỉnh, và được tổng kết, đánh giá trong chương kết luận tiếp theo.
